\section{The welfare ranking and European policy}
\label{sec:synthesis}\label{sec:policy}

The ex-post instruments of Section~\ref{sec:shield} and the ex-ante tax of
Section~\ref{sec:exante} can now be placed on one welfare scale, and the ranking
read against what conditions it.

\paragraph{The welfare ranking.}\label{sec:ranking} Table~\ref{tab:summary_price_import} reports welfare, measured as the total
consumption-equivalent variation (CEV) relative to the common no-tax steady
state. On private welfare the Slutsky transfer dominates the cap under
\emph{both} shocks (price: $-0.88\%$ vs.\ $-1.15\%$; supply: $-0.24\%$ vs.\ $-0.71\%$); the case for the transfer over the cap
does not depend on the transition. But the \emph{direction} of the cap's welfare
effect depends on the supply elasticity: tightening the cap raises welfare under
the price shock and lowers it under the supply shock (Figure~\ref{fig:dose}).

\input{output/tab_summary_import.tex}

\paragraph{Reading the ranking.} Three results point the same way. Energy-indexed targeting
is dominated by a flat transfer of equal budget (Section~\ref{sec:flat}); a
permanent carbon tax is mildly welfare-improving in the steady state through a
non-Pigouvian adoption correction, but essentially worthless as crisis insurance
(Sections~\ref{sec:carbontax}--\ref{sec:versionA}); and its insurance value is zero
to first order and does not grow as shocks become persistent, because greening
ex ante exhausts the adoption absorber (Section~\ref{sec:routeA}). The margin this
paper adds is a stabiliser, not an instrument to activate before the crisis: it
pays off when the economy is left free to adopt in the crisis.

\paragraph{Private welfare versus the transition.} The CEV above is a
\emph{private} household-welfare measure: it values consumption and leisure along
the transition and prices no climate or transition externality. Two results are
therefore kept separate. On private welfare alone the transfer dominates the cap
under both shocks; the case for the transfer does not depend on the transition.
Separately, the transfer preserves the green adoption margin the cap eliminates
(Section~\ref{sec:cap}): a positive statement about the transition rather than a
welfare ranking. We do not aggregate the two into one social-welfare number: that would
require a social value of adoption (a carbon price times avoided brown-energy use)
we do not model. A positive such value only reinforces the private ranking; it
cannot reverse it, since the transfer already dominates.

\paragraph{The supply elasticity and the Langot--Bayer contrast.}\label{sec:supply_elasticity} Under the supply shock the cap yields cumulative output of $-45.1$ against
$-14.9$ with no policy: with inelastic supply, capping the price prevents
substitution away from energy, so the economy cannot adjust its energy use at all.
This result survives the addition of an adoption margin.

\citet{langot2026tariff},
assuming elastic supply, find the French shield effective; \citet{bayer2026hank},
assuming fixed union-wide supply, find subsidies destructive. Our two closures
reproduce the \emph{sign} each obtains, which locates one source of the contrast
in a single parameter. Once the supply elasticity is fixed, the direction of the
cap's welfare effect is determined, so part of the empirical disagreement reduces
to measuring that elasticity at the relevant horizon.

The results map directly into the European policy architecture, on both the
ex-post and the ex-ante side.

\paragraph{Shield incomes rather than prices.} The \euro{600}-billion-plus shield of
2021--23 was dominated by price-suppressing measures, caps, regulated tariffs,
energy tax cuts \citep{Bruegel2023}. In the model, price suppression delivers output stabilisation but sacrifices the
entire crisis-time adoption wave: holding the brown retail price fixed removes the
price signal that rewards switching, and the extensive margin shuts
(Section~\ref{sec:shield}). The
data support the mechanism. European heat pump sales, having surged 38\% to a
record three million units in 2022 when gas was expensive, fell 5\% in 2023 and
21\% in 2024 as gas prices normalised, and the industry association's own account
names ``the low price of subsidised gas'' among the three main causes
\citep{EHPA2024,EHPA2025}. A transfer of the same fiscal cost stabilises at
least as well and leaves the wave intact. The 2022 episode left behind a fiscal bill and a forfeited adoption wave.

\paragraph{Pay the Social Climate Fund flat.} The EU's carbon price for
buildings and road transport (ETS2), scheduled to start in 2028 after a
one-year postponement, comes attached to a Social Climate Fund that will
mobilise at least \euro{86.7} billion over 2026--2032, of which up to 37.5\%
may finance temporary direct income support to vulnerable households
\citep{ECETS2}. Our distributional result speaks to the design of that component:
because energy use tracks total consumption within technology groups, indexing
support to energy expenditure targets wealth, not need. A flat payment of the
same cost dominates the energy-indexed transfer in welfare, and the gain is
concentrated on the constrained households the Fund is meant to
protect (Section~\ref{sec:shield}). Income support that preserves the price
signal is the model's best ex-post instrument; income support indexed to
fossil-energy use is a weaker version of it.

\paragraph{Defend ETS2 as externality pricing rather than insurance.} A recurring
argument for pre-crisis carbon pricing is resilience: a greener economy, it is
said, is better insured against the next fossil shock. The model disciplines
that claim. The crisis dividend of the prepared economy is zero to first order,
shrinks with the size and persistence of the shock, and the mechanism is
structural: pre-greening spends the adoption absorber, so the prepared economy
enters the crisis with less welfare-relevant room to green, even though it
switches more in volume (Section~\ref{sec:exante}). The case for ETS2 is the externality and the standing welfare gain of the greener
steady state (in the model a permanent brown tax raises standing welfare
monotonically over the range we can solve, by about $+0.28\%$ at a $10\%$ rate,
through a non-Pigouvian adoption correction); it is not protection against the
next crisis. The insurance rationale is the one the model supports least.
